\section{One periodic support and its integer secant remainder}\label{sec:secant}

Let an orbit of least period $n$ have successive states
$P_i=(x_i,x_{i+1})$, with indices read modulo $n$. Its coordinate
word satisfies
\begin{equation}\label{eq:recurrence}
 p(x_i)=x_{i+1}+a x_{i-1}.
\end{equation}
Let $E=\{x_i:0\le i<n\}$, put $m=\min E$ and
$D=\max E-m$. The integer translation
$(x,y)=(m+u,m+v)$ replaces $p$ by
\begin{equation}\label{eq:translation}
 p_m(t)=p(m+t)-(a+1)m\in\ZZ[t].
\end{equation}
It preserves degree, the Jacobian sign, the integer lattice and
native least period. We may therefore assume
$\min E=0$ and $\max E=D$. Unlike a global point-count argument,
these extrema are taken for one orbit, not for all periodic points
of a map. If $D=0$, the only state is $(0,0)$ and the period is one.
Assume below that $D\ge1$, and again denote the translated polynomial
by $p$.

For each actual letter $t\in E$, equation~\eqref{eq:recurrence} gives
\begin{equation}\label{eq:outputs}
 p(t)\in E+aE\subset J_a(D),\qquad
 J_1(D)=[0,2D],\quad J_{-1}(D)=[-D,D].
\end{equation}
Both intervals have length $2D$. Set
\begin{equation}\label{eq:secantparams}
 A=p(0),\qquad q=\frac{p(D)-p(0)}D.
\end{equation}
An integral-coefficient polynomial satisfies
$s-t\mid p(s)-p(t)$ for distinct integers $s,t$.
Thus $q\in\ZZ$; the endpoint output bounds imply $-2\le q\le2$.

\begin{lemma}[Integer secant remainder]\label{lem:secant}
There exists $h\in\ZZ[t]$ such that
\begin{equation}\label{eq:secant}
 p(t)=A+qt+t(t-D)h(t).
\end{equation}
For each actual interior letter $t\in E\cap(0,D)$,
\begin{equation}\label{eq:errorbound}
 |h(t)|\,t(D-t)\le2D.
\end{equation}
For all distinct integers $s,t$ one also has
\begin{equation}\label{eq:congruence}
 h(s)-h(t)\in(s-t)\ZZ.
\end{equation}
\end{lemma}
\begin{proof}
The polynomial $p(t)-A-qt$ is integral and vanishes at $0$ and $D$.
Divide first by the monic factor $t$. The quotient remains integral
and vanishes at $D$, since $D\ne0$; division by $t-D$ gives
\eqref{eq:secant}. No assumption on the leading coefficient was used.
The line $A+qt$ interpolates the two endpoint outputs, so its value
lies in $J_a(D)$ throughout $[0,D]$. At an actual letter, both it
and $p(t)$ lie in that interval of length $2D$. Their difference
proves \eqref{eq:errorbound}. Formula~\eqref{eq:congruence} is the
usual value-difference divisibility for an integral polynomial.
\end{proof}

If $D\ge10$, the bound in Lemma~\ref{lem:secant} gives
\begin{equation}\label{eq:endbounds}
 \begin{array}{c|c}
 \text{actual letter }t&\text{possible integer }h(t)\\ \hline
 1,\ D-1&-2,-1,0,1,2\\
 2,\ D-2&-1,0,1\\
 3\le t\le D-3&0
 \end{array}
\end{equation}
The last row follows from $t(D-t)\ge3(D-3)>2D$; it remains strict
at $D=10$. These constraints are imposed only on letters belonging
to $E$. Missing integers between the endpoints are not inserted into
the orbit.
